\documentclass[10pt,a4paper]{amsart}
\usepackage[margin=19mm]{geometry}
\usepackage{amsmath,amssymb,mathtools}
\usepackage[scr=rsfs,cal=euler]{mathalfa}
\usepackage{xfrac}
\usepackage[unicode,hidelinks,bookmarks=false]{hyperref}
\newcommand{\ie}{i.e.\ }
\newcommand{\cf}{cf.\ }
\newcommand{\resp}{respectively\ }
\newcommand{\Subsection}{\subsection}
\providecommand{\nobreakdash}{\nobreak\hbox{-}}
\setlength{\emergencystretch}{2em}
\multlinegap0pt
\usepackage{bm}
\usepackage{algorithm}\usepackage{algpseudocode}
\usepackage{xcolor}
\newtheorem{assumption}{Assumption}
\newtheorem{lemma}{Lemma}
\newtheorem{theorem}{Theorem}
\title{Distributed MPC For Coordinated Path-Following}
\author{Lusine Poghosyan, Anna Manucharyan, Mikayel Aramyan, Naira Hovakimyan, Tigran Bakaryan}
\date{Source: arXiv:2603.24748v1 (CC BY 4.0)}
\begin{document}
\maketitle

\noindent\emph{Source and licence.} Distributed MPC For Coordinated Path-Following. Version arXiv:2603.24748v1 (CC BY 4.0), distributed by arXiv under the Creative Commons Attribution 4.0 International licence, http://creativecommons.org/licenses/by/4.0/, as recorded in the arXiv licence metadata for this exact version and shown on the abstract page https://arxiv.org/abs/2603.24748v1. Full licence notice: \texttt{licenses/arxiv-2603.24748-CC-BY-4.0.txt}.\par\smallskip

\noindent\textit{Editorial scope.} Counted unit: the article's theorem theorem-unconstraint (source0.tex lines 822--841). The article's complete original proof of it is delivered with it (source0.tex lines 842--1169, one \texttt{proof} environment of 328 lines). No mathematics is written, repaired or re-derived: the statement and the proof are the article's own lines, in the article's own notation, and no retained line is substituted or re-flowed.  Retained uncounted as the source-local context the proof needs: lines 387--393; lines 396--403; lines 572--578; lines 663--679; lines 696--705; lines 710--717.  Nothing was omitted from any retained span, so the package carries no omission record.  No retained line contains a citation, so the wrapper carries no bibliography and no bibliography entry.  No retained line contains a figure environment, an image-inclusion command or drawing code, so no image is delivered and the package declares no asset dependency.  Editorial formatting only, in addition: an AMS \texttt{amsart} preamble that restates the article's own declarations and macros used by the retained lines, namely the environments \texttt{assumption}, \texttt{lemma}, \texttt{theorem} on separate counters, together with the standard packages those lines need.  The retained environments print in this document as Assumption 1, Lemma 1, Theorem 1 in order of appearance, whereas the article prints the same items with the numbering of its own full document.  The complete native source of the article as distributed in the arXiv e-print for this exact version is bundled unchanged as \texttt{sources/197167/source0.tex}.
\begin{assumption}\label{ass:graph}
The communication topology is fixed, i.e.,
\[
\mathcal N_{i,k}=\mathcal N_i,\qquad i=1,\dots,N,\quad k\ge 0,
\]
and the associated communication graph is undirected and connected.
\end{assumption}

\begin{assumption}\label{ass:tracking}
The path-following correction satisfies
\[
|\alpha_i^k|\le d e^{-\nu t_k}=d e^{-\nu kh},
\qquad i=1,\dots,N,\quad k\ge 0,
\]
for some constants \(d>0\) and \(\nu>0\).
\end{assumption}

\begin{equation}\label{varsub}
\bar{\bm{\Delta}}^k = V^{-1}\bm{\Delta}^k,
\quad
\dot{\bar{\bm{\Delta}}}^k = V^{-1}\dot{\bm{\Delta}}^k,
\quad
\bar{\bm{\alpha}}^k = V^{-1}\bm{\alpha}^k .
\end{equation}

\begin{equation}\label{state5}
\begin{bmatrix}
\bar{\delta}_{i,1}^k \\[5pt]
\dot{\bar{\delta}}_{i,1}^k
\end{bmatrix}
=
(Q^{h}_{i})^k
\begin{bmatrix}
\bar{\delta}_{i,1}^{0} \\[5pt]
\dot{\bar{\delta}}_{i,1}^{0}
\end{bmatrix}
+
\begin{bmatrix}
\bar{v}_i^k \\[5pt]
\bar{w}_i^k
\end{bmatrix},
\end{equation}

\begin{equation}\label{vk_def}
\begin{bmatrix}
\bar{v}_i^k \\[5pt]
\bar{w}_i^k
\end{bmatrix}
=
\sum_{s=1}^{k}
(Q_i^h)^{k-s}\, B^h\, \bar{\alpha}_i^{s},
\qquad k=1,2,\dots .
\end{equation}

\begin{lemma}\label{lemma:rho}
For every eigenvalue \(\lambda_i>0\) of \(D^{-1}L\), there exists \(h_i>0\) such that, for all
\(h\in(0,h_i)\),
\[
\rho(Q_i^h)<1,
\]
where \(\rho(Q_i^h)\) denotes the spectral radius of \(Q_i^h\).
\end{lemma}

\begin{theorem}\label{theorem-unconstraint}
Suppose that Assumptions~\ref{ass:graph}--\ref{ass:tracking} hold. Then, for sufficiently small \(h>0\), there exist
\[
r_h\in(0,1),\qquad M_0^h>0,\qquad \dot M_0^h>0,
\]
such that, for all \(k\ge 1\),
\begin{equation}\label{eq:teo1}
    |\delta_{i,0}^k-\delta_{j,0}^k|
\le M_0^h r_h^k,
\,\ |\dot{\delta}_{i,0}^{k}|
\le \dot M_0^h r_h^k,\,\ i,j=1,\dots,N.
\end{equation}
In particular,
\[
\delta_{i,0}^k-\delta_{j,0}^k\to 0,
\qquad
\dot{\delta}_{i,0}^k\to 0,
\quad \text{as } k\to\infty.
\]
\end{theorem}

\begin{proof}
We first establish convergence results for the auxiliary variables defined in \eqref{varsub}.
In particular, we distinguish two cases based on the spectrum of \(D^{-1}L\): strictly positive eigenvalues and the zero eigenvalue. For the first case, Lemma~\ref{lemma:rho} implies exponential convergence of the corresponding variables to zero. For the second case, we show convergence to a constant.
Combining these results and mapping them back to the original variables concludes the proof.

Since the communication graph is undirected and connected, the eigenvalues of \(D^{-1}L\) satisfy
\[
0=\lambda_1<\lambda_2\le \dots \le \lambda_N\le 2.
\]
Hence, for \(i=2,\dots,N\), one has \(\lambda_i>0\), and Lemma~\ref{lemma:rho} implies that, for sufficiently small \(h>0\),
\[
\rho(Q_i^h)<1,\qquad i=2,\dots,N.
\]
Fix such an \(h>0\), and choose \(r_h\in(0,1)\) such that
\[
\rho(Q_i^h)<r_h<1,\,\, i=2,\dots,N,
\,\,
r_h>e^{-\nu h},
\,\,
r_h>1-hb^h.
\]
By Gelfand's formula,
\[
\lim_{k\to\infty}\|(Q_i^h)^k\|_\infty^{1/k}
=
\rho(Q_i^h)
<r_h,
\qquad i=2,\dots,N.
\]
Therefore, for each \(i=2,\dots,N\), there exists a constant \(C_{i,\infty}^h>0\) such that
\[
\|(Q_i^h)^k\|_\infty \le C_{i,\infty}^h\, r_h^k,
\qquad \forall k\in\mathbb N.
\]
Since the set \(\{2,\dots,N\}\) is finite, defining
\[
C_\infty^h:=\max_{i=2,\dots,N} C_{i,\infty}^h,
\]
we obtain
\[
\|(Q_i^h)^k\|_\infty \le C_\infty^h\, r_h^k,
\qquad \forall k\in\mathbb N,\quad i=2,\dots,N.
\]

For \(i=2,\dots,N\), relation \eqref{state5} gives
\[
\begin{bmatrix}
\bar{\delta}_{i,1}^k \\[5pt]
\dot{\bar{\delta}}_{i,1}^k
\end{bmatrix}
=
(Q_i^h)^k
\begin{bmatrix}
\bar{\delta}_{i,1}^{0} \\[5pt]
\dot{\bar{\delta}}_{i,1}^{0}
\end{bmatrix}
+
\begin{bmatrix}
\bar{v}_i^k \\[5pt]
\bar{w}_i^k
\end{bmatrix}.
\]
Using the bound above, \eqref{vk_def}, submultiplicativity of induced norms, we obtain
\[
\left\|
\begin{bmatrix}
\bar{\delta}_{i,1}^k \\[5pt]
\dot{\bar{\delta}}_{i,1}^k
\end{bmatrix}
\right\|_\infty
\le
C_\infty^h r_h^k
\left\|
\begin{bmatrix}
\bar{\delta}_{i,1}^{0} \\[5pt]
\dot{\bar{\delta}}_{i,1}^{0}
\end{bmatrix}
\right\|_\infty
+
C_\infty^h\|B^h\|_\infty
\sum_{s=1}^{k}
r_h^{k-s}|\bar{\alpha}_i^s|.
\]
Moreover, by  Assumption~\ref{ass:tracking} for \(i=1,\dots,N\),
\[
|\bar{\alpha}_i^s|
\le \|\bar{\bm{\alpha}}^s\|_\infty
\le \|V^{-1}\|_\infty \|\bm{\alpha}^s\|_\infty
\le \|V^{-1}\|_\infty d\,e^{-\nu sh}.
\]
Therefore,
\[
\begin{split}
\left\|
\begin{bmatrix}
\bar{\delta}_{i,1}^k \\[5pt]
\dot{\bar{\delta}}_{i,1}^k
\end{bmatrix}
\right\|_\infty
&\le
C_\infty^h r_h^k \|V^{-1}\|_\infty
\max\{\|\bm{\Delta}^0\|_\infty,\|\dot{\bm{\Delta}}^0\|_\infty\}
\\
&\quad+
C_\infty^h\|B^h\|_\infty \|V^{-1}\|_\infty d
\sum_{s=1}^{k}
r_h^{k-s}e^{-\nu sh}.
\end{split}
\]
Since \(r_h>e^{-\nu h}\),
\[
\sum_{s=1}^{k}
r_h^{k-s}e^{-\nu sh}
=
r_h^k\sum_{s=1}^{k}\left(\frac{e^{-\nu h}}{r_h}\right)^s
\le
r_h^k\frac{e^{-\nu h}}{r_h-e^{-\nu h}}.
\]
Hence, for all \(i=2,\dots,N\) and all \(k\ge 0\),
\begin{equation}\label{eq:disagreement_mode_bound}
\left\|
\begin{bmatrix}
\bar{\delta}_{i,1}^k \\[5pt]
\dot{\bar{\delta}}_{i,1}^k
\end{bmatrix}
\right\|_\infty
\le
A_1^h
\max\{\|\bm{\Delta}^0\|_\infty,\|\dot{\bm{\Delta}}^0\|_\infty\}\,r_h^k
+
A_2^h\, d\, r_h^k,
\end{equation}
where
\[
A_1^h:=C_\infty^h\|V^{-1}\|_\infty,
\,\,
A_2^h:=C_\infty^h\|B^h\|_\infty \|V^{-1}\|_\infty
\frac{e^{-\nu h}}{r_h-e^{-\nu h}}.
\]

We now consider the case \(i=1\), corresponding to \(\lambda_1=0\). In this case,
\[
Q_1^h=
\begin{bmatrix}
1 & h-\dfrac{h^2}{2}b^h\\[6pt]
0 & 1-hb^h
\end{bmatrix},
\]
so \(Q_1^h\in\mathbb R^{2\times2}\) has two distinct eigenvalues
\[
\mu_1=1,
\qquad
\mu_2=q:=1-hb^h,
\qquad 0<q<1.
\]
Its characteristic polynomial is
\[
\chi(x)=(x-1)(x-q).
\]
By the Cayley--Hamilton theorem,
\[
(Q_1^h)^2-(1+q)Q_1^h+qI=0.
\]
Hence, for every \(k\in\mathbb N\), \((Q_1^h)^k\) can be written in the form
\[
(Q_1^h)^k=c_k Q_1^h+d_k I.
\]
Imposing this identity at the eigenvalues \(1\) and \(q\) gives
\[
c_k+d_k=1,
\qquad
c_k q+d_k=q^k,
\]
and therefore
\[
(Q_1^h)^k
=
\frac{1}{1-q}(Q_1^h-qI)+\frac{q^k}{1-q}(I-Q_1^h).
\]
Setting
\[
P:=\frac{1}{1-q}(Q_1^h-qI),
\qquad
F:=\frac{1}{1-q}(I-Q_1^h),
\]
we have
\begin{equation}\label{eq:Q_power_decomp_thm_refined}
(Q_1^h)^k=P+q^kF.
\end{equation}
Because $0<q<1$, we deduce that
\[
\lim_{k\to\infty}(Q_1^h)^k=P.
\]
Since the second row of \(P\) is zero, \eqref{eq:Q_power_decomp_thm_refined} implies
\[
\begin{split}
&\left|
\left[
(Q_1^h)^k
\begin{bmatrix}
\bar{\delta}_{1,1}^{0}\\[5pt]
\dot{\bar{\delta}}_{1,1}^{0}
\end{bmatrix}
\right]_2
\right|
\le
q^k\|F\|_\infty
\left\|
\begin{bmatrix}
\bar{\delta}_{1,1}^{0}\\[5pt]
\dot{\bar{\delta}}_{1,1}^{0}
\end{bmatrix}
\right\|_\infty
\\&\quad\le
r_h^k\|F\|_\infty\|V^{-1}\|_\infty
\max\{\|\bm{\Delta}^0\|_\infty,\|\dot{\bm{\Delta}}^0\|_\infty\}.
\end{split}
\]
Also, substituting \eqref{eq:Q_power_decomp_thm_refined} into \eqref{vk_def}, we obtain
\[
\begin{bmatrix}
\bar{v}_1^k \\[5pt]
\bar{w}_1^k
\end{bmatrix}
=
PB^h\sum_{s=1}^{k}\bar{\alpha}_1^s
+
FB^h\sum_{s=1}^{k}q^{k-s}\bar{\alpha}_1^s.
\]
The first term converges as \(k\to\infty\), since \(\sum_{s=1}^\infty \bar{\alpha}_1^s\) converges absolutely, while the second component of the first term is identically zero. Moreover,
\[
\sum_{s=1}^{k} q^{k-s}|\bar{\alpha}_1^s|
\le
\|V^{-1}\|_\infty d
\sum_{s=1}^{k} q^{k-s}e^{-\nu sh}
\le C_q^h r_h^k
\]
for some constant \(C_q^h>0\), because \(q<r_h\) and \(e^{-\nu h}<r_h\). Therefore, there exists \(\widetilde M^h>0\) such that
\[
|\dot{\bar{\delta}}_{1,1}^{k}|
\le
\widetilde M^h r_h^k.
\]

Returning to the original coordinates, since
\[
\dot{\bm{\Delta}}^k = V\dot{\bar{\bm{\Delta}}}^k,
\]
we obtain
\[
|\dot{\delta}_{i,1}^{k}|
\le
\|V\|_\infty \max_{s=1,\dots,N}|\dot{\bar{\delta}}_{s,1}^{k}|,
\qquad i=1,\dots,N.
\]
Using \eqref{eq:disagreement_mode_bound} for \(s=2,\dots,N\) and the previous estimate for \(s=1\), we conclude that there exists \(\dot M^h>0\) such that
\begin{equation}\label{eq:rate_bound_index1}
|\dot{\delta}_{i,1}^{k}|
\le
\dot M^h r_h^k,
\qquad i=1,\dots,N.
\end{equation}

Next, since \(\bar{\delta}_{1,1}^k\) is common to all agents in the reconstruction of \(\bm{\Delta}^k\), and the first column of \(V\) is \(\mathbf 1=[1,\dots,1]^T\), we have
\[
\delta_{i,1}^k-\delta_{j,1}^k
=
\sum_{s=2}^{N}(v_{i,s}-v_{j,s})\,\bar{\delta}_{s,1}^{k},
\qquad i,j=1,\dots,N.
\]
Hence, by \eqref{eq:disagreement_mode_bound},
\[
\begin{split}
|\delta_{i,1}^k-\delta_{j,1}^k|
&\le
\sum_{s=2}^{N}|v_{i,s}-v_{j,s}|\,|\bar{\delta}_{s,1}^{k}|
\\
&\le
2\|V\|_\infty \max_{s=2,\dots,N} |\bar{\delta}_{s,1}^{k}|
\\
&\le
M_1^h r_h^k,
\qquad i,j=1,\dots,N,
\end{split}
\]
for some constant \(M_1^h>0\).

Finally, using the shift conditions
\[
\delta_{i,0}^k=\delta_{i,1}^{k-1}-\alpha_i^{k},
\qquad
\dot\delta_{i,0}^k=\dot\delta_{i,1}^{k-1},
\qquad k\ge 1,
\]
we obtain
\[
|\dot\delta_{i,0}^k|
=
|\dot\delta_{i,1}^{k-1}|
\le
\dot M^h r_h^{k-1}
=
\frac{\dot M^h}{r_h}\,r_h^k,
\qquad i=1,\dots,N,
\]
and
\[
\begin{split}
|\delta_{i,0}^k-\delta_{j,0}^k|
&\le
|\delta_{i,1}^{k-1}-\delta_{j,1}^{k-1}|+|\alpha_i^k|+|\alpha_j^k|
\\
&\le
M_1^h r_h^{k-1}+2d\,e^{-\nu hk}
\\
&\le
\left(\frac{M_1^h}{r_h}+2d\right) r_h^k,
\qquad i,j=1,\dots,N,
\end{split}
\]
since \(e^{-\nu h}<r_h\). Therefore, the claim follows with
\[
\dot M_0^h:=\frac{\dot M^h}{r_h},
\qquad
M_0^h:=\frac{M_1^h}{r_h}+2d.
\]
This completes the proof.
\end{proof}

\end{document}
